\documentclass[11pt,letterpaper]{article}
\usepackage[utf8]{inputenc}
\usepackage[spanish,es-tabla]{babel}
\usepackage[margin=2.5cm]{geometry}
\usepackage{amsmath,amssymb}
\usepackage{booktabs}
\usepackage{graphicx}
\usepackage{float}
\usepackage{fancyhdr}
\usepackage{hyperref}
\usepackage[table]{xcolor}
\usepackage{tikz}
\usepackage{transparent}
\usepackage{eso-pic}
\usepackage{caption}
\definecolor{GeoOrange}{HTML}{FF8C00}
\definecolor{GeoBlue}{HTML}{0D47A1}
\AddToShipoutPictureBG{
  \AtPageCenter{
    \makebox[0pt]{\transparent{0.06}
\includegraphics[width=12cm]{Iconos/logo_geopol.png}}
  }
}
\hypersetup{colorlinks=true, linkcolor=GeoBlue, urlcolor=GeoOrange}
\pagestyle{fancy}
\fancyhf{}
\fancyhead[L]{\footnotesize \textcolor{GeoBlue}{\textbf{GeoPol Web}} - Reporte Técnico}
\fancyhead[R]{\footnotesize Universidad Distrital F.J.C.}
\fancyfoot[C]{\thepage}
\renewcommand{\headrulewidth}{0.4pt}
\renewcommand{\footrulewidth}{0.4pt}
\begin{document}
\begin{titlepage}
\begin{tikzpicture}[remember picture,overlay]
  \fill[GeoBlue] (current page.north west) rectangle ([yshift=-4cm]current page.north east);
  \fill[GeoOrange] ([yshift=-4cm]current page.north west) rectangle ([yshift=-4.5cm]current page.north east);
  \fill[GeoBlue!5] (current page.south west) rectangle ([yshift=4cm]current page.south east);
\end{tikzpicture}
\vspace*{-2cm}
\begin{center}
  \textcolor{white}{\Huge \textbf{GEOPORTAL WEB (GeoPol)}} \\[0.5cm]
  \textcolor{white}{\large \textit{''Máxima precisión al alcance de tus manos''}} \\[2cm]
  \vspace{3cm}
  {\LARGE \textbf{INFORME TÉCNICO DE INGENIERÍA}} \\[0.5cm]
  {\Large \textbf{Memorias de Cálculo Matemático y Diseño Vial\\ (Cubicaje de Volúmenes)}} \\[2cm]
  \begin{flushleft}
    \Large \textbf{Autores del Procesamiento:}\\[0.2cm]
    \large $\bullet$ Kevin Stiven Cubillos Ramirez \\[0.1cm]
    \large $\bullet$ Sergio Eduardo Barbosa Torres \\[0.1cm]
    \vspace{1cm}
    \Large \textbf{Tutor - Director de Proyecto:}\\[0.2cm]
    \large $\bullet$ Ing. Edgar Ladino
  \end{flushleft}
  \vfill
  \textbf{UNIVERSIDAD DISTRITAL FRANCISCO JOSÉ DE CALDAS}\\[0.2cm]
  Facultad de Medio Ambiente y Recursos Naturales \\ Ingeniería Topográfica / Civil \\[0.2cm]
  Bogotá D.C. -- 28 de July de 2026
\end{center}
\end{titlepage}
\tableofcontents
\newpage
\section{Introducción y Metodología de Cubicaje}
Este reporte detalla el cálculo de movimiento de tierras para el proyecto vial. El cálculo de volúmenes se ejecutó mediante el \textbf{Método de las Áreas Medias (Average End Area Method)}, evaluando sección transversal por sección transversal de forma secuencial.
\subsection{Criterios del Diagrama de Masas (Curva Masa)}
Para el análisis del movimiento de tierras, se ha calculado el volumen neto acumulado. La Curva Masa permite visualizar analíticamente el déficit o superávit de material a lo largo del eje del proyecto, facilitando el diseño de las zonas de préstamo y botadero.
\section{Resumen Ejecutivo de Volúmenes}
\begin{itemize}
  \item \textbf{Volumen Total de Corte (Excavación):} 401.280 $m^3$
  \item \textbf{Volumen Total de Relleno (Terraplén):} 72.480 $m^3$
  \item \textbf{Balance Neto del Proyecto:} 328.800 $m^3$
\end{itemize}
\section{Cuadro de Movimiento de Tierras}
A continuación, se presenta la tabla de cálculo tabulada. Los valores negativos indican predominancia de relleno, mientras que los positivos corresponden a áreas de corte.
\begin{table}[H]
  \centering
  \small
  \caption{Cuadro generalizado de cubicaje (Parte 1)}
  \label{tab:cubicaje_0}
  \begin{tabular}{|c|c|c|c|c|}
    \hline
    \rowcolor{GeoBlue}
    \textcolor{white}{\textbf{Abscisa (K)}} & \textcolor{white}{\textbf{Área Corte (m²)}} & \textcolor{white}{\textbf{Área Relleno (m²)}} & \textcolor{white}{\textbf{Vol. Corte (m³)}} & \textcolor{white}{\textbf{Vol. Relleno (m³)}} \\
    \hline
    \rowcolor{GeoBlue!5}
    0.000 & 24.720 & 0.000 & 0.000 & 0.000 \\
    \hline
    \rowcolor{white}
    10.000 & 9.723 & 0.003 & 172.215 & 0.015 \\
    \hline
    \rowcolor{GeoBlue!5}
    20.000 & 0.795 & 6.075 & 52.590 & 30.390 \\
    \hline
    \rowcolor{white}
    30.000 & 15.720 & 0.000 & 82.575 & 30.375 \\
    \hline
    \rowcolor{GeoBlue!5}
    40.000 & 3.060 & 2.340 & 93.900 & 11.700 \\
    \hline
  \end{tabular}
\end{table}
\vspace{0.8cm}
\begin{table}[H]
  \centering
  \small
  \caption{Cuadro generalizado de cubicaje (Parte 2)}
  \label{tab:cubicaje_1}
  \begin{tabular}{|c|c|c|c|c|}
    \hline
    \rowcolor{GeoBlue}
    \textcolor{white}{\textbf{Abscisa (K)}} & \textcolor{white}{\textbf{Vol. Acumulado Corte}} & \textcolor{white}{\textbf{Vol. Acumulado Relleno}} & \textcolor{white}{\textbf{Volumen Neto (m³)}} & \textcolor{white}{\textbf{Masa Acumulada (m³)}} \\
    \hline
    \rowcolor{GeoBlue!5}
    0.000 & 0.000 & 0.000 & 0.000 & 0.000 \\
    \hline
    \rowcolor{white}
    10.000 & 172.215 & 0.015 & 172.200 & 172.200 \\
    \hline
    \rowcolor{GeoBlue!5}
    20.000 & 224.805 & 30.405 & 194.400 & 366.600 \\
    \hline
    \rowcolor{white}
    30.000 & 307.380 & 60.780 & 246.600 & 613.200 \\
    \hline
    \rowcolor{GeoBlue!5}
    40.000 & 401.280 & 72.480 & 328.800 & 942.000 \\
    \hline
  \end{tabular}
\end{table}
\vspace{0.8cm}

\section{Diagrama de Masas (Curva Masa)}
Evolución gráfica del volumen acumulado en función de la abscisa (Distancia en K).
\begin{figure}[H]
  \centering
  \includegraphics[width=0.95\textwidth]{Reportes_PDF/Curva_Masa.png}
  \\caption{Diagrama de Masas para compensación longitudinal de material}
\end{figure}
\section{Conclusiones y Dictamen Técnico}
\begin{itemize}
  \item Dado el balance volumétrico resultante, el proyecto requiere de un transporte de material excedentario hacia un sitio de disposición final (botadero), ya que los volúmenes de excavación superan a los terraplenes requeridos.
  \item El \textbf{Diagrama de Masas} demuestra los puntos críticos de corte y la distribución longitudinal del material, permitiendo a los ingenieros y planificadores viales organizar el acarreo en volquetas de manera eficiente.
\end{itemize}
\newpage
\section{Anexo Gráfico: Perfiles de Secciones Transversales}
A continuación, se adjuntan los perfiles de todas las abscisas calculadas. Conforme a las normativas de presentación de planos de diseño vial, las secciones se han ploteado ordenadas de forma ascendente: \textbf{de abajo hacia arriba y de izquierda a derecha}.
\begin{figure}[H]
  \centering
  \begin{minipage}{0.48\textwidth}
    \vspace{4.5cm}
  \end{minipage}
  \hfill
  \begin{minipage}{0.48\textwidth}
    \vspace{4.5cm}
  \end{minipage}
  \\[0.3cm]
  \begin{minipage}{0.48\textwidth}
    \includegraphics[width=\linewidth]{Reportes_PDF/Sec_K40.000.png}
  \end{minipage}
  \hfill
  \begin{minipage}{0.48\textwidth}
    \vspace{4.5cm}
  \end{minipage}
  \\[0.3cm]
  \begin{minipage}{0.48\textwidth}
    \includegraphics[width=\linewidth]{Reportes_PDF/Sec_K20.000.png}
  \end{minipage}
  \hfill
  \begin{minipage}{0.48\textwidth}
    \includegraphics[width=\linewidth]{Reportes_PDF/Sec_K30.000.png}
  \end{minipage}
  \\[0.3cm]
  \begin{minipage}{0.48\textwidth}
    \includegraphics[width=\linewidth]{Reportes_PDF/Sec_K0.000.png}
  \end{minipage}
  \hfill
  \begin{minipage}{0.48\textwidth}
    \includegraphics[width=\linewidth]{Reportes_PDF/Sec_K10.000.png}
  \end{minipage}
  \\[0.3cm]
  \caption{Anexo Secciones Transversales - Plancha 1}
\end{figure}
\end{document}